% Bagean: metodhe
% Bab: bukti
% Pérangan: bukti-kanthi-kontradiksi

\documentclass[../../../include/open-logic-section]{subfiles}

\begin{document}

\olfileid{mth}{prf}{con}

\olsection{Bukti Kanthi Kontradiksi}

Wiwitané, bukti kanthi kontradiksi iku pola inferensi kanggo
mbuktekaké pratelan négatif. Anggepen panjenengan arep
nuduhaké yèn pratelan~$p$ iku \emph{salah}, yaiku
nuduhaké~$\lnot p$. Cara kang paling njanjèni yaiku
(a) nganggep $p$~bener, banjur (b) nuduhaké yèn anggepan
iki marakaké bab kang wis dingertèni salah.
``Bab kang dingertèni salah'' bisa awujud asil kang
nentang, utawa kontradiktif karo, $p$ dhéwé utawa
hipotesis liya saka pratelan kang lagi dirembug.
Upamané, bukti kanggo ``yèn $q$ mula $\lnot p$''
nganggep $q$~bener lan mbuktekaké~$\lnot p$ saka iku.
Yèn $\lnot p$ dibuktèkaké kanthi kontradiksi, kita
uga nganggep $p$ saliyané~$q$. Yèn kita
bisa mbuktekaké $\lnot q$ saka $p$, kita wis nuduhaké yèn
anggepan~$p$ marakaké bab kang nentang anggepan
liyane~$q$, amarga $q$ lan $\lnot q$ ora bisa padha
bener. Mesthi waé, mbuktekaké kontradiksi isih
mbutuhaké pola inferensi liya lan njlentrehaké
dhéfinisi. Ayo deleng sawijining conto.

\begin{prop}
  Yèn $A \subseteq B$ lan $B = \emptyset$, mula $A$
  ora nduwèni !!{element}s.
\end{prop}

\begin{proof}
  Anggepen $A \subseteq B$ lan $B = \emptyset$.
  Kita arep nuduhaké yèn $A$ ora nduwèni !!{element}s.
  \begin{quote}
    Amarga iki pratelan kondisional, kita nganggep
    antesedèn lan arep mbuktekaké konsekuèné.
    Konsekuèné yaiku: $A$ ora nduwèni !!{element}s.
    Iki bisa dipratelakaké luwih gamblang: ora bener
    yèn ana~$x \in A$.
  \end{quote}
  $A$ ora nduwèni !!{element}s yèn lan mung yèn
  ora bener ana~$x$ kang nyukupi $x \in A$.
  \begin{quote}
    Dadi pratelan kang arep dibuktèkaké iku pratelan
    négatif~$\lnot p$: ora bener yèn ana $x \in A$.
    Kanggo bukti kanthi kontradiksi, kita kudu nganggep
    pratelan positif kang cocog~$p$, yaiku ana
    $x \in A$, banjur nurunaké kontradiksi saka iku.
    Kita nandhani cara iki kanthi nulis
    ``kanthi anggepan kang ndadèkaké kontradiksi''
    utawa mung ``anggepen kosok baliné'', banjur
    nyebut anggepan~$p$.
  \end{quote}
  Anggepen kosok baliné: ana $x \in A$.
  \begin{quote}
    Iki anggepan anyar kanggo ngasilaké kontradiksi.
    Isih ana rong anggepan liya: $A \subseteq B$
    lan $B = \emptyset$. Anggepan kapisan mènèhi
    $x \in B$:
  \end{quote}
  Amarga $A \subseteq B$, $x \in B$.
  \begin{quote}
    Nanging amarga $B = \emptyset$, saben !!{element}
    saka $B$, upamané $x$, uga kudu !!a{element}
    saka~$\emptyset$.
  \end{quote}
  Amarga $B = \emptyset$, $x \in \emptyset$.
  Iki kontradiksi, amarga miturut dhéfinisi $\emptyset$
  ora nduwèni !!{element}s.
  \begin{quote}
    Sejatiné iki wis ngrampungaké bukti: saka
    anggepan-anggepan mau kita olèh kontradiksi,
    mula ora kabèh anggepan bisa bener bebarengan.
    Rong anggepan kapisan, $A \subseteq B$ lan
    $B = \emptyset$, tetep ditampa. Dadi anggepan
    kang pungkasan, yèn ana $x \in A$, mesthi salah.
    Nanging yèn kepéngin luwih rinci, iki bisa ditulis.
  \end{quote}
  Mula anggepan yèn ana $x \in A$ mesthi salah;
  kanthi bukti kontradiksi, $A$ ora nduwèni !!{element}s.
\end{proof}

Saben pratelan positif padha tegesé karo pratelan
négatif: $p$ yèn lan mung yèn $\lnot\lnot p$.
Dadi bukti kanthi kontradiksi uga bisa mbuktekaké
pratelan positif ``kanthi ora langsung''.
Kanggo mbuktekaké~$p$, wacanen minangka
pratelan négatif $\lnot\lnot p$. Yèn saka
$\lnot p$ kita bisa nurunaké kontradiksi,
mula $\lnot\lnot p$ wis kabuktèkaké kanthi
kontradiksi, lan mulané~$p$.

Ing conto mau kita arep mbuktekaké pratelan négatif,
yaiku $A$ ora nduwèni !!{element}s. Mula anggepan
kanggo bukti kontradiksi, yaiku ana $x \in A$,
iku pratelan positif. Pratelan iki mènèhi bahan
kanggo nalar, yaiku $x \in A$ kang hipotètis,
nganti tekan $x \in \emptyset$.

Yèn pratelan positif dibuktèkaké kanthi cara ora
langsung, anggepan kanggo bukti kontradiksi
awujud négatif. Nanging pratelan positif asring
gampang dipratelakaké minangka pratelan négatif,
lan kosok baliné. Bukti sadurungé sejatiné
padha yèn kita mbuktekaké ``$A = \emptyset$''
tinimbang konsekuèn négatif ``$A$ ora nduwèni
!!{element}s''. Miturut dhéfinisi $=$,
``$A = \emptyset$'' iku pratelan umum, amarga
tegesé ``saben !!{element} saka~$A$ uga
!!{element} saka~$\emptyset$ lan kosok baliné''.
Nanging iki cetha padha tegesé karo pratelan
négatif ``ora bener yèn ana $x \in A$''.

Mula kadhangkala luwih gampang nganggep $\lnot p$
tinimbang mbuktekaké~$p$ kanthi langsung.
Senajan bukti langsung padha gampang utawa
malah luwih gampang, kaya conto sabanjuré,
ana kang luwih seneng cara ora langsung.
Yèn negasi pindho mbingungaké, anggepen bukti
kontradiksi kanggo pratelan minangka cara
nurunaké kontradiksi saka pratelan kang
\emph{kosok baliné}. Dadi bukti kontradiksi
kanggo $\lnot p$ miwiti saka anggepan~$p$;
déné bukti kontradiksi kanggo~$p$ miwiti
saka~$\lnot p$.

\begin{prop}
$A \subseteq A \cup B$.
\end{prop}

\begin{proof}
  Kita arep nuduhaké yèn $A \subseteq A \cup B$.
  \begin{quote}
    Sekilas iki pratelan positif: saben $x \in A$
    uga ana ing $A \cup B$. Negasiné yaiku:
    ana $x \in A$ kang $x \notin A \cup B$.
    Mula pratelan iki bisa dibuktèkaké ora langsung
    kanthi nganggep negasiné lan nurunaké
    kontradiksi.
  \end{quote}
  Anggepen kosok baliné, yaiku
  $A \nsubseteq A \cup B$.
  \begin{quote}
    Dhéfinisi $A \subseteq A \cup B$ yaiku:
    saben $x \in A$ uga $x \in A \cup B$.
    Kanggo ngerti tegesé $A \nsubseteq A \cup B$,
    dhéfinisi kang wis dijlentrehaké kudu
    dianalisis kanthi logika dhasar: pratelan
    yèn saben $x \in A$ uga $x \in
    A \cup B$ iku salah yèn lan mung yèn ana
    \emph{sawijining}~$x \in A$ kang
    $x \notin A \cup B$.
    Ing kéné kudu ngati-ati: akèh bukti kontradiksi
    kang dicoba déning mahasiswa gagal amarga
    salah nganalisis negasi pratelan kaya
    ``kabèh $A$s iku $B$s''. Tegesé,
    $A \nsubseteq A \cup B$ yèn lan mung yèn
    ana $x$ kang nyukupi $x \in A$ lan
    $x \notin A \cup B$. Sabanjuré gampang.
  \end{quote}
  Dadi ana $x \in A$ kang nyukupi
  $x \notin A \cup B$. Miturut dhéfinisi
  $\cup$, $x \in A \cup B$ yèn lan mung yèn
  $x \in A$ utawa $x \in
  B$. Amarga $x \in A$, kita olèh
  $x \in A \cup B$. Iki nentang anggepan
  $x \notin A \cup B$.
\end{proof}

\begin{prob}
Buktèkna \emph{kanthi ora langsung} yèn
$A \cap B \subseteq A$.
\end{prob}

\begin{prop}
Yèn $A \subseteq B$ lan $B \subseteq C$,
mula $A \subseteq C$.
\end{prop}

\begin{proof}
  Anggepen $A \subseteq B$ lan $B \subseteq C$.
  Kita arep nuduhaké $A
  \subseteq C$.
  \begin{quote}
    Kita gunakna cara ora langsung: anggepen
    negasi saka kang arep dibuktèkaké.
  \end{quote}
  Anggepen kosok baliné, yaiku
  $A \nsubseteq C$.
  \begin{quote}
    Kaya mau, $A \nsubseteq C$ yèn lan mung
    yèn ora saben $x \in A$ uga $x \in C$,
    yaiku ana $x \in A$ kang $x \notin C$.
    Aja kuwatir; kanthi latihan, njlentrehaké
    negasi kaya iki bakal dadi lumrah.
  \end{quote}
  Tegesé, ana~$x$ kang nyukupi $x \in A$
  lan $x \notin C$.
  \begin{quote}
    Saiki iki bisa digunakaké kanggo tekan
    kontradiksi. Mesthi waé kita uga kudu
    migunakaké rong anggepan liyane.
  \end{quote}
  Amarga $A \subseteq B$, $x \in B$.
  Amarga $B \subseteq C$, $x \in
  C$. Nanging iki nentang $x \notin C$.
\end{proof}

\begin{prop}
Yèn $A \cup B = A \cap B$, mula $A = B$.
\end{prop}

\begin{proof}
  Anggepen $A \cup B = A \cap B$.
  Kita arep nuduhaké yèn $A = B$.
  \begin{quote}
    Wiwitané saiki wis lumrah:
  \end{quote}
  Anggepen, kanggo ngasilaké kontradiksi,
  yèn $A \neq B$.
  \begin{quote}
    Anggepan kanggo bukti kontradiksi yaiku
    $A \neq B$. Amarga $A = B$ yèn lan mung
    yèn $A \subseteq B$ lan $B \subseteq A$,
    kita olèh $A \neq B$ yèn lan mung yèn
    $A \nsubseteq B$ \emph{utawa}
    $B \nsubseteq
    A$. Élinga yèn negasi kudu diolah kanthi
    ngati-ati!{} Kanggo nurunaké kontradiksi
    saka disjungsi iki, kita mbédakaké kasus
    lan nuduhaké kontradiksi ing saben kasus.
  \end{quote}
  $A \neq B$ yèn lan mung yèn $A \nsubseteq B$
  utawa $B \nsubseteq A$. Kita mbédakaké kasus.
  \begin{quote}
    Ing kasus kapisan, anggepen $A \nsubseteq B$,
    yaiku ana $x$ kang $x \in A$ nanging
    $x \notin B$. $A \cap B$ yaiku
    himpunan !!{element}s kang padha-padha
    ana ing $A$ lan $B$. Mula yèn sawijining
    obyèk ora ana ing salah siji, obyèk iku
    uga ora ana ing irisan. $A \cup B$
    nggabungaké $A$ lan $B$, mula anggota
    salah sijiné uga anggota gabungan.
    Iki nuduhaké $x \in A \cup B$ nanging
    $x \notin A \cap
    B$, lan mula $A \cap B \neq A \cup B$.
  \end{quote}

  Kasus 1: $A \nsubseteq B$. Mula ana $x$,
  $x \in A$ nanging $x \notin
  B$. Amarga $x \notin B$, mula
  $x \notin A \cap B$. Amarga $x \in A$,
  $x \in A \cup B$. Dadi
  $A \cap B \neq A \cup B$, kang nentang
  anggepan $A \cap B = A \cup B$.

  Kasus 2: $B \nsubseteq A$. Mula ana $y$,
  $y \in B$ nanging $y \notin
  A$. Kaya mau, kita olèh $y \in A \cup B$
  nanging $y \notin A \cap B$, mula
  $A \cap B \neq A \cup B$, kang maneh
  nentang $A \cap B = A \cup
  B$.
\end{proof}

\end{document}
